% ── Section 3: Methodology ──────────────────────────────────────────────────
\section{Methodology}
\label{sec:methodology}

\subsection{Overview and Pipeline Flow}
The multimodal digital twin framework is structured as a hierarchical, multi-stage prognostic pipeline that mirrors clinical decision-making under real-world observational constraints (Fig.~\ref{fig:architecture}). Rather than operating as an isolated single-modality diagnostic calculator, the system integrates point-in-time cross-sectional staging, multi-horizon trajectory forecasting, sequential consultation history modeling, counterfactual scenario simulation, and interpretability governance into a unified workflow.

\begin{figure}[H]
    \centering
    \includegraphics[width=0.95\linewidth]{figures/fig_architecture.pdf}
    \caption{End-to-end system architecture: Multimodal patient intake routes clinical, cognitive, fluid, and imaging data through the calibrated TwinStager (Model A), the multi-horizon progression engine (Model B), and the Counterfactual Scenario Lab, bounded by empirical plausibility guards and TreeSHAP explainability.}
    \label{fig:architecture}
\end{figure}

The pipeline executes through four sequential operational stages:
\begin{enumerate}
    \item \textbf{Intake and Current Stage Detection (Model A)}: Heterogeneous patient observations across demographics, bedside cognitive scores, cerebrospinal fluid (CSF) assays, and structural neuroimaging are ingested without statistical imputation. The calibrated multi-tier TwinStager dynamically identifies the highest clinically valid feature tier and computes calibrated posterior probabilities across Cognitively Normal (CN), Mild Cognitive Impairment (MCI), and Alzheimer's Disease (AD).
    \item \textbf{Longitudinal Progression Tracking (Model B)}: Anchor visit cognitive scores and predicted staging probabilities feed into target-specific delta-change regressors that forecast prospective decline across 6-, 12-, and 24-month horizons for three core cognitive endpoints: the Mini-Mental State Examination (MMSE), the Clinical Dementia Rating Global score (CDGLOBAL), and the Alzheimer's Disease Assessment Scale--Cognitive Subscale (ADAS11-Cog).
    \item \textbf{Conversion Risk Estimation}: For baseline non-demented anchor visits (CN and MCI), binary hazard classifiers estimate the probability of progression to AD across 6, 12, and 24 months. Projected test scores are subsequently re-routed through Model~A to estimate complete future stage probability distributions at each horizon.
    \item \textbf{Sequence-Based Staging (Bidirectional LSTM)}: When a patient presents with multiple historical consultations, a Bidirectional Long Short-Term Memory (BiLSTM) network processes chronological sliding windows ($T \le 3$) to capture temporal momentum and predict syndromic transitions at the subsequent visit.
    \item \textbf{Counterfactual Simulation and Governance}: The Counterfactual Scenario Lab receives projected trajectories, allowing clinicians to perturb cognitive decline velocity or simulate literature-grounded lifestyle interventions. Every projection is audited against an empirical 5th--95th percentile training envelope, physiological boundary guards, TreeSHAP feature attributions, leave-one-modality-out sensitivity shifts, and an explicit 43-field intake lineage matrix.
\end{enumerate}

\subsection{Data, Cohort, and Leakage Controls}
The primary longitudinal dataset was constructed from the Alzheimer's Disease Neuroimaging Initiative (ADNI) repository~\cite{petersen2010adni}, consolidating participants across ADNI-1, ADNI-GO, ADNI-2, and ADNI-3. A master dataset comprising \TotalRows{} clinical visits across \TotalSubjects{} unique participants was compiled by integrating five primary diagnostic modalities: clinical diagnosis consensus records, bedside neuropsychological examination batteries, cerebrospinal fluid molecular assays, ComBat-harmonized FreeSurfer volumetric MRI, and APOE targeted genotyping (Table~\ref{tab:cohort_splits} and Table~\ref{tab:data_sources}).

\input{tables/table1_cohort_splits.tex}
\input{tables/table2_data_sources.tex}

To guarantee statistical validity and prevent observational bias, four architectural anti-leakage safeguards were enforced (Fig.~\ref{fig:data_pipeline}):
\begin{enumerate}
    \item \textbf{Subject-Level Stratified Partitioning}: Rather than partitioning records at the visit level, which allows longitudinal visits from the same subject to appear in both training and evaluation partitions, data splitting was performed strictly at the subject level. Partitioning utilized a stratified group split conditioned on participant baseline diagnosis and longitudinal conversion status \texttt{(baseline\_diagnosis, ever\_converted)}. As summarized in Table~\ref{tab:cohort_splits}, \TrainSubjects{} subjects (\TrainPct{}) were allocated to model training, \ValSubjects{} subjects (\ValPct{}) to hyperparameter validation, and \TestSubjects{} subjects (\TestPct{}) to an independent held-out evaluation test set. Converter proportions were exactly matched across partitions (17.5\%), and programmatic assertions verified zero subject overlap between splits.
    \item \textbf{Strict Backward-Only Neuroimaging Temporal Joins}: In routine neurology, neuroimaging is rarely obtained on the exact date of clinical examination. To align imaging with cognitive evaluations without leaking future anatomical deterioration, each cognitive visit was linked exclusively to the most recent volumetric MRI scan acquired on or prior to the visit date. Forward-looking scans (acquired after the examination date) were strictly excluded. Scans acquired more than 180 days prior were treated as unobserved missing entries, preventing temporally decoupled brain morphometry from biasing current staging.
    \item \textbf{Visit Timeline Reconstruction}: Exact clinical examination dates were extracted from recorded visit timestamps (72.0\%) or calculated from recorded baseline anchor dates and protocol visit intervals (28.0\%). Baseline and screening visits occurring within protocol proximity were collapsed to visit month $0.0$, eliminating duplicate baseline anchor points.
    \item \textbf{Forbidden Feature Blacklisting}: Sixteen features encoding full longitudinal history or diagnostic targets were permanently blacklisted. These include total lifetime visit counts, retrospective longitudinal trajectory classifications, imputation confidence flags, and retrospective clinician diagnosis revision codes. Models trained on such variables exhibit artificially inflated test accuracy by inadvertently reading future patient outcomes.
    \item \textbf{Handling Observational Missingness Without Imputation}: In real-world cohorts, diagnostic missingness is non-random; invasive tests such as lumbar punctures or research-grade MRI are administered selectively based on clinical suspicion. Imputing unobserved values via mean, median, or multivariate imputation (MICE) distorts physiological variance and leaks cross-subject future information. We strictly rejected statistical imputation: missing modalities (84.7\% missingness for CSF, 45.5\% for volumetric MRI) were retained as native missing values and routed through tree-split decision logic.
\end{enumerate}

\begin{figure}[H]
    \centering
    \includegraphics[width=0.95\linewidth]{figures/fig_data_pipeline.pdf}
    \caption{Longitudinal data pipeline and leakage controls: (Left) Strictly backward-only temporal merge enforcing zero future scan attachment ($\le 180$ days); (Right) Subject-level stratified partition preserving converter proportions across splits.}
    \label{fig:data_pipeline}
\end{figure}

\subsection{Step 1, Current Stage Detection: Multi-Tier Stager (Model A)}
Cross-sectional syndromic staging classifies each clinical consultation into one of three standard clinical categories: Cognitively Normal (CN), Mild Cognitive Impairment (MCI), or Alzheimer's Disease (AD). Because clinical centers operate under varying diagnostic resource constraints—ranging from primary care clinics with bedside screening to specialized memory centers with fluid biomarkers and volumetric MRI—Model~A was designed with a multi-tier feature routing architecture. Rather than requiring complete multimodal profiles, the model dynamically evaluates the patient's record and routes inference through the highest valid tier:
\begin{itemize}
    \item \textbf{Tier T1 (Demographics \& Genetics)}: Patient entry age, biological sex, completed education years, and APOE $\varepsilon4$ risk allele count ($0, 1, 2$).
    \item \textbf{Tier T1b}: Tier T1 augmented with the Mini-Mental State Examination (\texttt{MMSE}), representing rapid bedside cognitive screening.
    \item \textbf{Tier T2a}: Tier T1b augmented with the Clinical Dementia Rating Global score (\texttt{CDGLOBAL}) and Sum of Boxes (\texttt{CDRSB}), capturing functional impairment.
    \item \textbf{Tier T2b (Core Cognitive Battery)}: Tier T2a augmented with detailed neuropsychological batteries: ADAS11-Cog total score (\texttt{TOTSCORE}), the Functional Activities Questionnaire (\texttt{FAQTOTAL}), and the Montreal Cognitive Assessment (\texttt{MOCA}), capturing fine-grained memory, praxis, and executive function.
    \item \textbf{Tier T3 (+CSF Biomarkers)}: Tier T2b augmented with cerebrospinal fluid amyloid-beta ($\text{A}\beta_{42}$), total tau (t-tau), hyperphosphorylated tau ($\text{p-tau}_{181}$), and lower/upper assay censorship flags.
    \item \textbf{Tier T4 (+FreeSurfer Volumetric MRI)}: Tier T3 augmented with ComBat-harmonized FreeSurfer morphometry, including normalized bilateral hippocampal volume, hippocampal asymmetry, ventricular volume, total brain parenchymal volume, and regional cortical thicknesses.
\end{itemize}

The stager is built using a CatBoostClassifier~\cite{prokhorenkova2018catboost}. CatBoost was selected over alternative architectures (e.g., Random Forests, Multi-Layer Perceptrons) because of its native handling of missing numerical values via optimal threshold routing during tree construction, robust handling of heterogeneous tabular distributions, resistance to overfitting via symmetric oblivious decision trees, and millisecond-scale deterministic inference. The stager was trained with 600 iterations, a maximum tree depth of 6, a learning rate of 0.05, multiclass cross-entropy loss, and balanced class weights ($1.0$ for CN, $1.5$ for MCI, and $2.0$ for AD) to account for lower relative prevalence of dementia in community cohorts. Early stopping was monitored on the validation partition with a 50-round patience window. Sample weights inversely proportional to subject visit counts ($w_i = 1 / N_{\text{visits}}$) were applied to prevent over-representation of subjects with extensive follow-up.

Because raw tree ensemble logits frequently exhibit uncalibrated overconfidence, we performed post-hoc calibration via Temperature Scaling—the multi-class generalization of Platt scaling. A scalar temperature parameter $T > 0$ was fitted on the validation partition by minimizing multiclass negative log-likelihood:
\begin{equation}
    P(y = k \mid \mathbf{x}) = \frac{\exp(z_k / T)}{\sum_{j=1}^{K} \exp(z_j / T)}
\end{equation}
where $z_k$ represents the uncalibrated raw margin logit for class $k$. Bounded scalar optimization yielded an optimal temperature of $T = 0.8554$, reducing the overall Expected Calibration Error (ECE) to $\StagerECEFinal{}$. In the main Results (Table~\ref{tab:stager_results}), empirical multi-tier evaluations demonstrate that diagnostic performance scales steadily from Tier T1 (Macro-F1 = \StagerTOneFOne{}) to Tier T4 (Macro-F1 = \StagerTFourFOne{}).

\subsection{Step 2, Progression Tracking: Multi-Horizon Delta Regression (Model B)}
Forecasting future cognitive decline is challenging because predicting absolute test scores ($\hat{y}_{t+h}$) directly often leads to compounding baseline drift. To anchor predictions to the patient's current disease state, Model~B is formulated as a delta-change regression model:
\begin{equation}
    \Delta y_{t, h} = y_{t+h} - y_t, \quad h \in \{6, 12, 24\text{ months}\}
\end{equation}
The future predicted score is reconstructed by adding the estimated change to the baseline anchor value: $\hat{y}_{t+h} = y_t + \widehat{\Delta y}_{t, h}$.

We selected three primary cognitive and functional endpoints:
\begin{enumerate}
    \item \textbf{MMSE} ($0$--$30$ scale): The primary clinical benchmark for global cognitive tracking;
    \item \textbf{CDGLOBAL} ($0$--$3$ scale): The standard ordinal measure of functional independence;
    \item \textbf{ADAS11-Cog} ($0$--$70$ scale): The primary cognitive endpoint in Alzheimer's disease pharmacological clinical trials.
\end{enumerate}
Progression horizons of 6, 12, and 24 months were selected to align with standard clinical follow-up schedules and therapeutic trial endpoints. For each endpoint and horizon, an independent CatBoostRegressor was trained on all valid anchor visits possessing an observed evaluation at horizon $h$. Features include baseline cognitive scores, age, sex, education, APOE $\varepsilon4$ allele count, and optional neuroimaging features.

To assess performance, Model~B was evaluated against two established prognostic baselines:
\begin{itemize}
    \item \textbf{Naive Persistence Baseline}: Assumes zero future change ($\widehat{\Delta y}_{t, h} = 0$, $\hat{y}_{t+h} = y_t$), representing the hypothesis that patient status remains static;
    \item \textbf{Linear Slope Extrapolation}: Computes the patient's historical annualized rate of change across prior visits ($\ge 2$ visits) and extrapolates linearly across horizon $h$.
\end{itemize}

Prediction intervals (90\% and 95\% confidence bounds) were generated using 1,000 cluster-bootstrap iterations resampled at the subject level. Resampling entire subject clusters with replacement accounts for repeated-measures intra-subject correlation, preventing the artificially narrowed confidence intervals produced by naive sample-level bootstrapping. As reported in Section~\ref{sec:results} (Table~\ref{tab:forecast_mae}), Model~B achieves an MAE of \MMSEMAEXII{} points for 12-month MMSE decline, significantly outperforming persistence (\MMSEPersistXII{}) and slope extrapolation (\MMSESlopeXII{}).

\subsection{Step 3, Conversion Risk Estimation}
A vital clinical task is predicting which non-demented patients will progress to dementia. Conversion risk estimation is restricted strictly to non-demented baseline anchor visits (Cognitively Normal and Mild Cognitive Impairment). Patients already diagnosed with Alzheimer's disease at the anchor visit are ineligible for conversion modeling, as they have already reached the dementia endpoint; for these patients, the progression engine instead evaluates rapid functional decline risk (probability of losing $\ge 3$ MMSE points).

For each horizon $h \in \{6, 12, 24\}$ months, a binary CatBoostClassifier was trained to predict whether an initial CN or MCI anchor will convert to AD within that window. Models were evaluated using two complementary metrics:
\begin{enumerate}
    \item \textbf{Area Under the Receiver Operating Characteristic (AUROC)}: Quantifies threshold-independent discrimination between converters and stable non-converters;
    \item \textbf{Brier Skill Score (BSS)}: Evaluates probabilistic calibration and refinement relative to a baseline reference forecast ($\text{BSS} = 1 - \text{BS} / \text{BS}_{\text{ref}}$). A positive BSS demonstrates that model probabilities outperform historical empirical incidence rates.
\end{enumerate}

To complete the longitudinal trajectory projection, projected cognitive scores ($\hat{y}_{t+h}$) are passed back into the calibrated TwinStager (Model A). This produces re-staged diagnostic posterior probabilities ($P(\text{CN})$, $P(\text{MCI})$, $P(\text{AD})$) at each future horizon, allowing clinicians to inspect expected syndromic transition timelines alongside continuous numerical decline (Section~\ref{sec:results}, Table~\ref{tab:conversion_metrics}, reaching AUROC of \ConvAUCXXIV{} and BSS of \ConvBSSXXIV{} at 24 months).

\subsection{Step 4, Sequence-Based Staging: Bidirectional LSTM}
While Models A and B operate on static single-visit anchor profiles, longitudinal consultation histories contain chronological trajectory signals. To capture sequence dynamics across multi-visit consultations, we incorporated a Bidirectional Long Short-Term Memory (BiLSTM) network. The model operates on sliding windows of up to $T = 3$ consecutive clinical consultations, mapping past score trajectories to next-visit syndromic stage.

Each consultation in the input sequence $\mathbf{X} \in \mathbb{R}^{B \times T \times D}$ tracks $D = 5$ standardized features: patient entry age, CDR global score (\texttt{CDGLOBAL}), Mini-Mental State Examination (\texttt{MMSCORE}), ADAS-Cog-11 prorated score (\texttt{TOTSCORE}), and elapsed visit month. Sequences with fewer than three visits are pre-padded with zeros along the temporal dimension ($T=3$). Features are normalized using an empirical Min-Max scaler fitted exclusively on the training partition. The network architecture consists of:
\begin{enumerate}
    \item A Bidirectional LSTM layer with 64 forward and 64 backward recurrent units ($128$ concatenated recurrent channels, input dropout $0.2$, recurrent dropout $0.1$) to capture both past deficit accumulation and retrospective rate-of-change;
    \item A Batch Normalization layer stabilizing hidden state activations across variable-length sequence inputs;
    \item A fully-connected Dense layer with 32 rectified linear units (ReLU) and $L_2$ kernel weight regularization ($\lambda = 10^{-4}$);
    \item A Dropout regularization layer ($p = 0.3$);
    \item A 3-unit Softmax classification layer computing posterior probabilities for CN, MCI, and AD at the subsequent clinical visit.
\end{enumerate}
The BiLSTM was trained on ADNI consultation sequences using the Adam optimizer ($\text{learning rate} = 10^{-3}$, batch size 32) and sparse categorical cross-entropy loss. Balanced inverse-frequency class weights were applied during training to counter the lower relative prevalence of dementia in early consultation visits.

\subsection{Recurrent Regression Experiment (Masked GRU)}
To investigate whether deep recurrent architectures could replace gradient boosted decision trees for continuous multi-horizon numerical regression across irregularly spaced visits, we conducted an empirical experiment with a two-layer Gated Recurrent Unit incorporating missingness indicator masks ($m_t \in \{0, 1\}$)~\cite{che2018grud}. The model was designed to handle observational visit sparsity by concatenating observation masks and time intervals into hidden state transitions.

A pre-defined architectural adoption criterion required the deep recurrent model to achieve a Macro-F1 improvement $\ge 0.03$ and lower MMSE MAE at 12 months on the held-out test partition before replacing CatBoost. While the masked GRU achieved an MMSE 12-month test MAE of \GRUMaskedMAE{} (significantly outperforming an unmasked recurrent GRU baseline of \GRUUnmaskedMAE{}, paired $\Delta = \GRUPairedDelta{}$, $p < 0.05$), CatBoost achieved a superior MAE of \GBMProgressionMAE{} points while delivering deterministic, sub-5ms CPU inference without requiring GPU infrastructure. Consequently, CatBoost was retained for the operational progression engine.

\subsection{Counterfactual Scenario Lab, Plausibility Envelope, and Lifestyle Modifier}
The Counterfactual Scenario Lab provides clinicians with an interactive simulation sandbox to test hypothetical risk modifications without altering the patient's official electronic record. Clinicians can adjust baseline cognitive scores or modulate expected decline velocity using an acceleration multiplier ($\alpha \in [0.5, 1.5]$).

To protect against implausible mathematical extrapolations, all scenario forecasts are bounded by three layers of clinical safeguards:
\begin{enumerate}
    \item \textbf{Empirical Training Envelope}: Derived from the 5th and 95th percentiles of the ADNI training cohort across anchor variables (Table~\ref{tab:archetype_plausibility}): age $[60.7, 84.7]$, MMSE $[20.0, 30.0]$, CDGLOBAL $[0.0, 1.0]$, and ADAS11 $[2.3, 26.3]$. When patient inputs fall outside these empirical boundaries, automated scope alerts are generated and prediction interval widths are expanded by $1.5\times$ to reflect elevated uncertainty.
    \item \textbf{Physiological Score Caps}: Predictions are bounded strictly to valid clinical assessment scales ($\text{MMSE} \in [0, 30]$, $\text{CDGLOBAL} \in [0, 3]$, $\text{ADAS11} \in [0, 70]$).
    \item \textbf{Biological Monotonicity Bounds}: Untreated neurodegenerative decline does not experience spontaneous physiological recovery. For patients staged as MCI or AD, positive MMSE changes are bounded to $\Delta \le +1.0$ point/year; for patients with advanced dementia ($\text{MMSE} \le 17$), positive score drift is restricted to zero.
\end{enumerate}

\input{tables/table8_archetype_plausibility.tex}

Because ADNI does not collect comprehensive longitudinal lifestyle logs, modifiable risk factors cannot be learned directly from cohort data without introducing synthetic bias. Instead, modifiable risks are incorporated via an external, literature-grounded prior derived from the Lancet Commission on Dementia Prevention, Intervention, and Care~\cite{livingston2020lancet,livingston2024lancet}. Ten modifiable risk factors are parameterized with pooled relative risks (RR) from international meta-analyses (Table~\ref{tab:lifestyle_factors}): current smoking ($\text{RR}=1.59$), physical inactivity ($\text{RR}=1.39$), excessive alcohol ($\text{RR}=1.19$), poor diet ($\text{RR}=1.20$), sleep disturbance ($\text{RR}=1.25$), midlife hypertension ($\text{RR}=1.61$), type 2 diabetes ($\text{RR}=1.52$), depression ($\text{RR}=1.91$), traumatic brain injury ($\text{RR}=1.78$), and cardiovascular disease ($\text{RR}=1.28$).

To prevent multiplicative compounding from inflating risk estimates when multiple comorbidities co-occur, relative risks are attenuated using square-root shrinkage ($\text{RR}_{\text{adj}} = \text{RR}^{\LifestyleShrinkage{}}$) and bounded by a conservative hard cap:
\begin{equation}
    M_{\text{lifestyle}} = \text{clip}\left(\prod_{i} \text{RR}_{i}^{\LifestyleShrinkage{}}, \, \LifestyleHardCapLo{}, \, \LifestyleHardCapHi{}\right)
\end{equation}
\textbf{Invariance Safeguard:} The lifestyle multiplier modulates prospective cognitive decline rates and conversion hazard probabilities only; it is mathematically barred from altering baseline syndromic staging probabilities.

\subsection{Explainability and Lineage Governance}
To ensure clinical transparency and support medical decision-making, the framework implements three interpretability mechanisms:
\begin{enumerate}
    \item \textbf{TreeSHAP Feature Attributions}: Exact Shapley additive explanations~\cite{lundberg2017shap} compute feature-level marginal contributions to CatBoost log-odds margins. For each prediction, attributions decompose the model's decision into positive and negative directional forces pushing toward CN, MCI, or AD.
    \item \textbf{Leave-One-Modality-Out Sensitivity Analysis}: For multimodal intakes, advanced diagnostic modalities (CSF biomarkers, volumetric MRI, detailed neuropsychological batteries) are sequentially suppressed to evaluate counterfactual stage probability shifts $\Delta P(\text{Stage})$. This reveals to clinicians whether invasive lumbar punctures or expensive MRI scans alter the diagnostic outcome relative to non-invasive cognitive tests alone.
    \item \textbf{43-Field Intake Lineage Matrix}: Every clinical variable collected at intake is mapped to an explicit system lineage tag (\texttt{used\_in\_stage}, \texttt{used\_in\_forecast}, \texttt{used\_in\_lifestyle}, \texttt{used\_in\_symptom}, \texttt{used\_in\_mri\_fusion}, \texttt{recorded\_only}) (Table~\ref{tab:lineage_summary}). This audit trail ensures that every variable's role is traceable, preventing unauthorized feature leakage between models.
\end{enumerate}

\subsection{Integration of Component Models and Contribution}
The system architecture coordinates the flow of variables and predictions across models:
\begin{itemize}
    \item Intake observations are routed to Model~A, which generates calibrated baseline stage probabilities;
    \item Baseline cognitive scores and Model~A stage probabilities pass to Model~B, which projects delta trajectories and conversion hazards;
    \item Multi-visit consultation histories are evaluated concurrently by the BiLSTM to provide sequential trajectory validation;
    \item Forecasted trajectories pass to the Counterfactual Scenario Lab, where literature-grounded lifestyle priors modulate decline rates within the empirical plausibility envelope;
    \item Interpretability modules extract TreeSHAP drivers, marginal modality shifts, and lineage audits across the entire flow.
\end{itemize}

The principal contributions of this integrated methodology are:
\begin{enumerate}
    \item A leak-free longitudinal data architecture enforcing subject-level partitioning, backward-only imaging joins ($\le 180$ days), and feature blacklisting without statistical imputation;
    \item A calibrated multi-tier routing stager that adapts dynamically to arbitrary clinical data availability;
    \item Open and transparent paired statistical evaluation demonstrating the lack of incremental utility for invasive CSF biomarkers and volumetric MRI;
    \item An audited counterfactual simulation environment bounded by empirical training percentiles and literature-grounded priors.
\end{enumerate}
